\subsection{Discussion}

Our results qualitatively support the central claim of \citet{darcet2024vision}: adding register tokens to a frozen ViT does not degrade downstream performance and slightly improves all three probes for DINOv2-L/14, the only cleanly comparable row. For OpenCLIP and DeiT-III, the absence of publicly released trained-register checkpoints forced us to use approximate substitutes (test-time registers and inference-time injection respectively), so those comparisons are weaker but still consistent with the claim. The Figure~8 ablation is the most difficult to reproduce: since retraining is out of scope, all three of our approaches (truncate/cycle, test-time registers, independent variant) test the effect of register tokens on a model that was not optimised for the chosen $N$, producing a near-binary rather than smooth curve. Within this constraint, the test-time register method of \citet{jiang2025vision} gives the most faithful proxy: one register absorbs the outlier and the metrics plateau, matching the paper's key observation. The unexpected $N=16$ collapse, which the independent variant confirms is not caused by value duplication, is a new finding beyond the scope of the original paper: it suggests that the Jiang method is correctly scoped at $N=1$ and that downstream blocks not trained on many broadcast tokens limit scalability.
